\documentclass[12pt,a4paper]{article}

\usepackage[utf8]{inputenc}
\usepackage[brazilian]{babel}
\usepackage{amsmath}
\usepackage{graphicx}
\usepackage{siunitx}
\usepackage{booktabs}
\usepackage{geometry}
\usepackage[authoryear]{natbib}

\geometry{margin=2.5cm}

\title{Determinação da Aceleração da Gravidade por Meio do Experimento de Queda Livre}
\author{Enzo Rocha Leite Diniz Ribas \\ Centro Federal de Educação Tecnológica de Minas Gerais}
\date{\today}

\begin{document}
\begin{titlepage}

\begin{center}


\includegraphics[width=4cm]{media/Logo_CEFET-MG.png}

\Large
Centro Federal de Educação Tecnológica de Minas Gerais

\vspace{1cm}

\Large
Curso de Engenharia Mecânica

\vspace{2cm}

\Huge
Determinação da Aceleração da Gravidade por Meio do Experimento de Queda Livre

\vspace{3cm}

\begin{flushright}
\Large
Mecânica, Oscilações, Fluidos e Termodinâmica, Turma 6T56 \\
Prof. Dr. Mauro Lucio Lobão Iannini
\end{flushright}

\vspace{3cm}

\Large
Enzo Rocha Leite Diniz Ribas (Eng. Mecânica - 20253002819)\\
Gabriel Resende Amorim Oliveira (Eng. Elétrica - 20233014438)\\
Leandro Césari e Silva Barros (Eng. Mecânica - 20253005104)\\

\vfill

\Large
Belo Horizonte

\Large
\the\year

\end{center}

\end{titlepage}

\section{Introdução}

O experimento de queda livre constitui um dos métodos clássicos para a determinação da aceleração da gravidade. A análise do movimento de um corpo submetido exclusivamente à ação da gravidade permite verificar experimentalmente as relações previstas pela cinemática do movimento uniformemente acelerado.

Neste trabalho busca-se determinar o valor local da aceleração da gravidade a partir da análise do tempo de queda de uma esfera metálica liberada a partir de diferentes alturas. Os dados experimentais foram analisados por meio de diferentes métodos de ajuste e linearização, permitindo comparar os resultados obtidos.

\section{Fundamentação Teórica}

Quando um objeto é abandonado a partir do repouso e submetido apenas à força gravitacional, ele realiza um movimento retilíneo uniformemente acelerado. Desprezando a resistência do ar, a equação do movimento é dada por

\begin{equation}
h = \frac{1}{2}gt^2
\end{equation}

onde:

\begin{itemize}
\item $h$ é a altura de queda,
\item $t$ é o tempo de queda,
\item $g$ é a aceleração da gravidade.
\end{itemize}

A partir dessa equação podem ser construídas diferentes representações gráficas para determinação experimental de $g$.

Uma possibilidade consiste em linearizar a equação construindo o gráfico

\begin{equation}
h \; \text{vs} \; t^2
\end{equation}

o que resulta em uma relação linear da forma

\begin{equation}
h = \left(\frac{g}{2}\right)t^2
\end{equation}

onde o coeficiente angular da reta corresponde a $g/2$.

Outra possibilidade consiste em isolar o tempo:

\begin{equation}
t = \sqrt{\frac{2}{g}}\sqrt{h}
\end{equation}

resultando em uma relação linear entre $t$ e $\sqrt{h}$.

\section{Metodologia}

O experimento foi realizado utilizando um sistema composto por sensores ópticos conectados a um cronômetro eletrônico. Uma esfera metálica foi liberada a partir de diferentes alturas e o tempo de queda foi registrado.

Foram selecionadas diferentes alturas ao longo da régua disponível, cobrindo uma faixa ampla de valores. Para cada altura, registrou-se o tempo correspondente à passagem da esfera entre os sensores.

Os dados foram organizados em tabela e analisados utilizando o software SciDAVis para geração de gráficos e realização de ajustes de curvas.

\section{Procedimentos}

O procedimento experimental seguiu as seguintes etapas:

\begin{itemize}
\item ajuste da altura inicial da esfera em diferentes posições da régua;
\item liberação da esfera a partir do repouso;
\item registro do tempo de queda utilizando sensores ópticos;
\item repetição das medições para diferentes alturas;
\item organização dos dados em tabela;
\item construção de gráficos e realização de ajustes no SciDAVis.
\end{itemize}

\section{Discussão e Análise dos Resultados}

A Tabela \ref{tab:dados} apresenta os valores experimentais obtidos.

\begin{table}[h]
\centering
\begin{tabular}{ccc}
\toprule
Altura $h$ (m) & Tempo $t$ (s) & $t^2$ (s$^2$) \\
\midrule
% inserir dados aqui
0.044 & 0.096 & 0.0092 \\
0.100 & 0.142 & 0.0202 \\
0.200 & 0.201 & 0.0404 \\
\bottomrule
\end{tabular}
\caption{Dados experimentais de altura e tempo de queda.}
\label{tab:dados}
\end{table}

\subsection{Gráfico $h$ vs $t$}

O gráfico da altura em função do tempo apresenta comportamento parabólico, conforme esperado teoricamente.

\begin{figure}[h]
\centering
\includegraphics[width=0.7\textwidth]{media/graph1_quadratic_fit.pdf}
\caption{Altura em função do tempo de queda.}
\end{figure}

\subsection{Linearização $h$ vs $t^2$}

A linearização dos dados permitiu obter uma relação aproximadamente linear.

\begin{figure}[h]
\centering
\includegraphics[width=0.7\textwidth]{media/graph1_quadratic_fit.pdf}
\caption{Gráfico linearizado $h$ em função de $t^2$.}
\end{figure}

A partir do ajuste linear obteve-se

\begin{equation}
g = 9.80 \pm 0.10 \; \text{m/s}^2
\end{equation}

\subsection{Linearização $t$ vs $\sqrt{h}$}

Também foi realizado um ajuste utilizando a relação

\begin{equation}
t = \sqrt{\frac{2}{g}}\sqrt{h} + b
\end{equation}

O ajuste forneceu

\begin{equation}
g = 10.55 \pm 0.14 \; \text{m/s}^2
\end{equation}

O valor de referência para a aceleração da gravidade em Belo Horizonte é

\begin{equation}
g_{BH} = 9.784 \; \text{m/s}^2
\end{equation}

O erro percentual foi calculado por

\begin{equation}
\text{Erro} = \frac{|g_{exp} - g_{ref}|}{g_{ref}} \times 100
\end{equation}

Diferenças entre os valores obtidos podem ser atribuídas a:

\begin{itemize}
\item incertezas na medição da altura;
\item atraso de resposta dos sensores;
\item pequenas variações no instante de liberação da esfera.
\end{itemize}

A resistência do ar foi considerada desprezível devido às pequenas alturas e às baixas velocidades envolvidas no experimento.

\section{Conclusão}

A análise experimental permitiu determinar a aceleração da gravidade com boa concordância com o valor esperado para a região de Belo Horizonte. Os diferentes métodos de ajuste apresentaram resultados próximos, sendo o método baseado na linearização $h$ vs $t^2$ o que apresentou maior estabilidade estatística.

\bibliographystyle{plainnat}
\bibliography{refs}

\end{document}